\begin{figure}[p]
\centering

\includegraphics[width=\linewidth]{a6efff7140e16fd202e094f504efbaea737deea9ccb83d4abcff42c748ab01a2.jpg}
\end{figure}

\begin{center}
{\large 香港中文大學}\\
The Chinese University of Hong Kong

\vspace{1em}

{\Huge 與自然對話}\\
{\Huge IN DIALOGUE WITH NATURE}

\vspace{1em}

通識教育基礎課程讀本\\
Textbook for General Education Foundation Programme

\vspace{1em}

Second edition, 2012
\end{center}

\begin{figure}[p]
\centering
\includegraphics[width=\linewidth]{8784b1a80b63978c6fc27b4602b394f6fdeabc934a2093867ec476486f1f88da.jpg}
\end{figure}

\section*{香港中文大學 大學通識教育部}
\begin{center}
Office of University General Education\\
The Chinese University of Hong Kong

\vspace{1em}

{\LARGE In Dialogue with Nature}\\
{\LARGE 與自然對話}
\end{center}

\noindent \textit{In Dialogue with Nature}（《與自然對話》）是香港中文大學通識教育基礎課程兩門課程之一「UGFN 1000 In Dialogue with Nature」的讀本。

\vspace{1em}

\noindent © Office of University General Education, The Chinese University of Hong Kong, 2011, 2012\\
Second edition 2012

\vspace{1em}

\noindent 版權所有。未經香港中文大學大學通識教育部事先書面許可，本出版物任何部分不得以任何形式複製，亦不得提供予他人使用。

\noindent 本書所選篇章均經版權持有人許可重印，或取自公有領域。除公有領域材料外，任何選文均不得在未取得「鳴謝」所列版權持有人事先許可的情況下複製。

\section*{Editorial Board}
Chan Chi Wang\\
Szeto Wai Man\\
Wong Wing Hung

\vspace{1em}

\noindent Managing editor / Ng Hiu Chun Emily\\
Editorial assistants / Lam Yee Ki Dora, Ng Tsz Ying Tracy\\
Cover and layout design / Ng Hiu Chun Emily

\vspace{1em}

\noindent ISBN: 978-988-16660-1-7

\vspace{1em}

\noindent Published by Office of University General Education\\
The Chinese University of Hong Kong\\
Shatin, Hong Kong

\vspace{1em}

\noindent http://www.cuhk.edu.hk/oge/gef\\
gef@cuhk.edu.hk

\vspace{1em}

\noindent Printed in Hong Kong by Wellfit Production \& Printing Company（唯美印刷製作公司）

\vspace{1em}

\noindent 本書以不經漂白和增白的原漿紙印刷。\\
This book is printed on unbleached, with no brightening agents, original pulp paper.

\section*{鳴謝}
\noindent 本書所選篇章均經版權持有人許可重印，或取自公有領域。除公有領域材料外，任何選文均不得在未取得下列版權持有人事先許可的情況下複製。

\vspace{0.5em}

\noindent pp. 5--10: From Plato, \textit{Republic}, translated by C.D.C. Reeve. Copyright © 2004 by Hackett Publishing Company, Inc. Reprinted by permission of Hackett Publishing Company, Inc. All rights reserved.

\noindent pp. 11--48: From David C. Lindberg, \textit{The Beginnings of Western Science: The European Scientific Tradition in Philosophical, Religious, and Institutional Context, Prehistory to A.D. 1450} (Second Edition). Published by The University of Chicago Press. Copyright © 1992, 2007 by The University of Chicago. All rights reserved.

\noindent pp. 49--62: From I. Bernard Cohen, \textit{The Birth of a New Physics: Revised and Updated}. Copyright © 1985 by I. Bernard Cohen. Copyright © 1960 by Educational Services Incorporated. Used by permission of W. W. Norton \& Company, Inc.

\noindent pp. 63--70: From Isaac Newton, \textit{The Principia: Mathematical Principles of Natural Philosophy}. Copyright © 2012 by University of California Press. Reproduced with permission of University of California Press via Copyright Clearance Center.

\noindent pp. 73--96: From Charles Darwin, \textit{On the Origin of Species by Means of Natural Selection} (First Edition) (1859). Reproduced from the public domain.

\noindent pp. 97--142: From James D. Watson and with Andrew Berry, \textit{DNA: The Secret of Life}. Copyright © 2003 by DNA Show LLC. Used by permission of Alfred A. Knopf, a division of Random House, Inc.

\noindent pp. 143--158: From Rachel Carson, \textit{Silent Spring: 40th Anniversary Edition}. Copyright © 1962 by Rachel L. Carson. Copyright renewed 1990 by Roger Christie. This selection originally appeared in \textit{SILENT SPRING: 40th Anniversary Edition} by Rachel Carson. This edition of ``pages 65 to pages 83 from \textit{SILENT SPRING: 40th Anniversary Edition} by Rachel Carson'' is published by Office of University General Education, The Chinese University of Hong Kong.

\noindent Published by arrangement with Frances Collin, Trustee u-w-o Rachel Carson through Bardon-Chinese Media Agency.

\noindent ALL RIGHTS RESERVED.

\noindent pp. 161--178: From Henri Poincaré, \textit{The Value of Science: Essential Writings of Henri Poincaré} (2001). Published by Modern Library, reproduced from the public domain.

\noindent pp. 179--194: From Eric Kandel, \textit{In Search of Memory: The Emergence of a New Science of Mind}. Copyright © 2006 by Eric R. Kandel. Used by permission of W. W. Norton \& Company, Inc.

\noindent pp. 195--218: From Joseph Needham and Colin A. Ronan, \textit{The Shorter Science and Civilisation in China: An Abridgement of Joseph Needham's Original Text}, Vol. 1 (1978). © Cambridge University Press, reproduced with permission.

\noindent pp. 219--244: Nathan Sivin, ``Why the Scientific Revolution Did Not Take Place in China -- or Didn't It?'' originally published in \textit{Chinese Science}, 1982, 5: 45--66, revised 2005.8.24, used by permission of the author.

\noindent pp. 245--252: 文言文篇錄自沈括，《新校正夢溪筆談》（1975），由中華書局（香港）有限公司授權翻印。詮譯及註釋由香港中文大學大學通識教育部提供。

\noindent pp. 253--258 (except for pp. 255--256 ``The Rainbow''): From Joseph Needham, \textit{Science and Civilisation in China}, Vol. 3 (1959), Vol. 4 (1962), Vol. 5 (1985). © Cambridge University Press, reproduced with permission.

\noindent pp. 255--256 (``The Rainbow''): Nathan Sivin, ``The Rainbow,'' unpublished, used by permission of the author.

\noindent pp. 259--274: From William Dunham, \textit{The Mathematical Universe: An Alphabetical Journey Through the Great Proofs, Problems, and Personalities}. Copyright © 1994 by John Wiley \& Sons, Inc. Reproduced with permission of John Wiley \& Sons, Inc.

\noindent pp. 275--290: From Euclid, \textit{Elements}. The text used in this anthology is adopted from a digitized version of Euclid, \textit{The Thirteen Books of Euclid's Elements}, translated by Thomas L. Heath (1956), made available by the Perseus Digital Library Project. Ed. Gregory R. Crane. Tufts University. (Date of access: Feb 2012) <http://www.perseus.tufts.edu>. The illustrations are reproduced by this publication.

\section*{目錄}
\noindent Letter to Students \dotfill vii\\
Notes on editing \dotfill xi\\
Introduction to ``In Dialogue with Nature'' \dotfill 1

\section*{Part I: Human Exploration of the Physical Universe}
\noindent Introduction \dotfill 3\\
Text 1a from \textit{Republic} / Plato \dotfill 5\\
Text 1b from \textit{The Beginnings of Western Science} / David C. Lindberg \dotfill 11\\
Text 2 from \textit{The Beginnings of Western Science} / David C. Lindberg \dotfill 17\\
Text 3a from \textit{The Birth of a New Physics} / I. Bernard Cohen \dotfill 49\\
Text 3b from \textit{The Principia: Mathematical Principles of Natural Philosophy} / Isaac Newton \dotfill 63

\section*{Part II: Human Exploration of the World of Life}
\noindent Introduction \dotfill 71\\
Text 4 from \textit{On the Origin of Species} / Charles Darwin \dotfill 73\\
Text 5 from \textit{DNA: The Secret of Life} / James D. Watson \dotfill 97\\
Text 6 from \textit{Silent Spring} / Rachel Carson \dotfill 143

\section*{Part III: Our Understanding of Human Understanding}
\noindent Introduction \dotfill 159\\
Text 7 from \textit{Science and Method} / Henri Poincaré \dotfill 161\\
Text 8 from \textit{In Search of Memory: The Emergence of a New Science of Mind} / Eric R. Kandel \dotfill 179\\
Text 9 from \textit{The Shorter Science and Civilisation in China} / Joseph Needham \dotfill 195\\
Text 10a from ``Why the Scientific Revolution Did Not Take Place in China -- or Didn't It?'' / Nathan Sivin \dotfill 219\\
Text 10b from 《夢溪筆談》 / 沈括 \dotfill 245\\
Brush Talks from Dream Brook / Shen Kua \dotfill 253\\
Text 11a from \textit{The Mathematical Universe} / William Dunham \dotfill 259\\
Text 11b from \textit{Elements} / Euclid \dotfill 275



\section*{致同學}
\noindent 親愛的同學：

\noindent 歡迎加入通識教育基礎課程（General Education Foundation Programme, GEF）的學習共同體。我們即將展開一段充滿挑戰的旅程：穿行於不同文化與學科的流水之間，進而航向遼闊的知識海洋。GEF 開設的兩門課程，即「與人文對話」（\textit{In Dialogue with Humanity}）與「與自然對話」（\textit{In Dialogue with Nature}），旨在拓闊視野，展示彼此不同的世界。我們將反思至關重要的議題與關懷，並在此過程中努力更完整地理解自身，以及我們所身處的宇宙。

\noindent 多數人進入大學，是為了在主修學科中獲得專門知識或一套技能，以此通向專業或學術生涯。然而，在當前這個技術快速變化、通訊高度整合、經濟波動不定的社會環境，也就是我們所稱的知識經濟之中，傳統的就業資歷已不再足夠。理解全球化世界的複雜性，需要寬廣而多元的視角。為了解決日益複雜的問題，往往必須同時而協調地調動多個知識領域。

\noindent 從根本上影響我們在世生活的，不只是我們擁有什麼知識或技能，也包括我們成為怎樣的人。因此，通識教育的終極目標，是培養負責任的知識分子與世界公民所需的素質。我們必須意識到認知世界的不同方式，回應人類存在的共同關切，對文化多樣性保持敏感；而最重要的是，我們必須培養面對陌生事物、並就其作出明智判斷的智性能力與心志。

\noindent 「與人文對話」與「與自然對話」構成香港中文大學通識教育課程的核心部分。這兩門課是以主題為本、由經典文本閱讀引導的研討課，也可從幾個方面被視為「核心」。

\noindent 就內容而言，GEF 課程涵蓋人類知識與實踐的兩個關鍵或核心領域：對作為人的意義與價值的追尋，以及人類理解自然的成就與限制。我們將閱讀的經典所包含的思想，觸及人類存在與知識中的重大問題，並界定了與當代生活仍然相關、歷久不衰的信念與價值。

\noindent 此外，兩門課也促進重要能力的發展，包括主動投入閱讀、討論、反思與寫作。這些能力既是終身學習者不可或缺的素質，也是大學生活以至往後人生得以充實展開所需的智性工具。

\noindent 更重要的是，GEF 課程提供了一個空間，讓核心態度得以開始形成。首先，是面對智性挑戰時的無畏精神。我們將直接與來自迥異傳統、文化與學科的原典對話，並努力理解它們。

\noindent 同樣重要的，是開放心靈的精神。當我們與文本、與彼此展開對話時，會發現各自的觀點常常互有差異。我們的信念與價值也可能受到挑戰。我們或會逐漸明白，對許多重要問題而言，並不存在確定或最終的答案；人類的知識與實踐是一個永無止境的尋求與再尋求的過程。也正是在這一過程中，我們將為負責任而積極地參與知識與福祉的追尋奠定基礎。

\noindent GEF 的參與者都是主動學習者。這一過程不應，也不會，被一個字母等級所概括。GEF 的學習成果包括：對文化差異更敏銳的覺察，以及對價值多樣性更深的敏感；對普遍關懷更強的意識，以及對人類知識的範圍與限制更深的敏感；面對艱深文本與複雜思想時的信心；以及以口頭或書面方式表達有根據論點的能力。要使這趟智性旅程成功，你們的主動參與正是關鍵。

\noindent 準備好吧！

\vspace{1em}

\noindent Leung Mei Yee\\
Director\\
General Education Foundation Programme

\section*{译者注}

\noindent 感谢强大的AI技术，让我可以毫无压力的完成对这本书的翻译。之前也有这本书的中文译本存在过，但那本书的翻译又臭又长还不准确。这本书略微的优化了翻译，但是也没好多少（笑）。本书完全由ChatGPT 5.5 进行翻译，译者本人没有进行过任何校对，但由于保留了英文原文，读者应该能很方便的进行对照和理解。希望这本书能对你的学习有所帮助。

\noindent 译者本人不对此书中的任何内容负责。

\noindent 另外，此译本完全免费发行，如果花钱你就被骗了。

\hfill 译者：智乃的脚小小的




\section*{編輯說明}
\begin{enumerate}
\item 本選集作為「與自然對話」課程的讀本出版；該課程是一個以主題為取向、以經典文本為本的核心課程的一部分。因此，對原典篇章的選取與刪節，均以使課程主題更為清晰集中為目標。
\item 本選集中重印的文本，除校正排印錯誤、為編輯一致性而調整排版樣式，以及刪去或調整原文本中對頁碼或附錄的引用外，均按原樣重印。
\item 書邊所有段落編號均由本書加上，以便課堂討論時引用。編號依照段落在原文本中出現的順序排列。因此，凡有段落被刪去，編號亦會跳過。
\item 所有腳註及方括號內文字均出自原文本；標明為「Ed.」者除外。
\item 原文本中的省略以方括號省略號「[...]」標示。凡如此標示的省略，均由本書作出。
\item 第 10b 篇「Brush Talks from Dream Brook 夢溪筆談」所有標題均由本書加上，並非原文本所有。中文文本中亦列出條目號碼，以標明相關條目在原文本中的出現次序。
\item 第 11b 篇「Elements」有以下特別說明：
\end{enumerate}

\begin{itemize}
\item 本選集只收入 Heath 的 \textit{Elements} 英譯，並刪去他對希臘文本的 ``scholia'' 評注以及譯註。
\item 在原出版物中，Heath 使用不同類型的括號，標示他為使譯文清楚而補入、但希臘原文中並非逐字具備的文字。為方便一般讀者閱讀，本書刪去這些括號而保留其中文字。唯一例外，是 Heath 用來指引讀者參照 \textit{Elements} 相關部分的方括號，例如「[Def. 10]」指定義 10，「[Post. 1]」指公設 1，「[C.N. 3]」指公理 3，「[I.4]」指第一卷命題 4；這些均予保留。
\end{itemize}

\section*{In Dialogue with Nature}
\section*{導論}
\noindent 人類長久以來都對自然懷有好奇。早在文明曙光出現以前，我們便一直努力理解世界所呈現的繁多現象。宇宙的某些運行似乎有序，例如日出日落、四季更替、天體旋轉；另一些則似乎混沌而隨機，例如風暴與火山爆發。或許，在這些模式或秩序背後，存在某些自然法則。為了理解這些觀察，人們在世界各地不同的具體境況與文化中，發展並確立了極為不同的理解自然的方法，並以之作為古代文化的基礎。

\noindent 早期自然哲學家及其後繼者，即今日的現代科學家，致力發現自然法則；這一努力一直是一場自然與我們之間的對話。這場對話透過經驗與推理進行。它始於約公元前 500 年的古希臘；當時哲學家開始把邏輯與理性的智性工具運用於對自然現象的觀察，以嘗試解釋存在。他們對自然的態度與探究方法，奠定了我們今日現代西方經驗科學及其所支撐技術的基礎。其後，更大量、更系統的觀察逐步修正了我們對自然的理解。現代科學家則從這些觀察中概括出自然定律。

\noindent 然而，在這一過程中，我們對自然的理解經歷了許多調整，甚至出現了一些極為深刻的根本變化。某些新的思維模式偶爾甚至挑戰了「理解」這一概念本身，迫使人類反思知識的意義。

\noindent 與此同時，在大致相同的時期，中國對理解自然的追尋走向了十分不同的答案。中國學者發展出陰陽、五行等概念，形成了一套與希臘同行迥然有別的理解。直到近代早期，西方科學的概念才開始滲入中國的文化精神；然而，中國在許多世紀中即使沒有被認為由科學帶來的那些益處，仍是古代世界公認的強國。顯然，東西文化之間持續不斷的相遇，對我們理解自然仍大有可為。

\noindent 本課程邀請同學追溯前人在這場求知歷程中的思路，並與他們的著作展開對話。沿著偉大心靈的足跡前行，同學將能對自然以及人類與自然的互動形成有根據的見解。

\section*{課程結構}
\noindent 本課程將帶領同學踏上一段智性旅程。我們首先走向外在世界，考察的對象是物質宇宙與生命世界。這兩項考察構成本課程的前兩部分：第一部分「人類對物質宇宙的探索」（Human Exploration of the Physical Universe）與第二部分「人類對生命世界的探索」（Human Exploration of the World of Life）。我們將認識西方科學若干重大發現：宇宙受物理定律支配；物種正在經歷自然選擇；DNA 是生命的密碼。

\noindent 接著，我們將回到自身，考察我們的心靈與理解。這是本課程第三部分「我們對人類理解的理解」（Our Understanding of Human Understanding），也是一段反思之旅。

\noindent 首先，我們將反思自己對自然的理解，以及科學知識的有效性與限制。我們將嘗試理解具備理性思考能力的人類心靈的本質。其次，透過比較西方科學與中國哲學，我們將嘗試理解為什麼現代科學沒有在中國出現；並且在中國哲學中尋找一些視角，幫助我們反思現代科學處理人與自然關係的方式。

\vspace{1em}

\noindent Wong Wing Hung\\
Deputy Director\\
General Education Foundation Programme

\section*{第一部分：人類對物質宇宙的探索}
\noindent 這是旅程的第一部分。我們將從柏拉圖的《理想國》（\textit{Republic}）開始。在通常被稱為「洞穴寓言」（Allegory of the Cave）的選段（Text 1a）中，柏拉圖透過蘇格拉底與格勞孔的對話，區分了感知與實在。這一區分奠定了西方科學的基本基調。

\noindent 依柏拉圖之見，我們在「可見領域」（visible realm）中所見的只是影子。只有在「可知領域」（intelligible realm）中，人才能找到真正實在的「善的理念」（Form of the Good），它是「一切正確而美好事物的原因」。這是真正追求理解者的終極關切。對現代科學家而言，也存在一個由自然法則構成的「可知領域」。科學家的終極關切，不是對自然現象作全面描述，而是理解其背後的法則。David Lindberg 是研究中世紀與近代早期科學的著名史家。在他廣受使用的教科書 \textit{The Beginnings of Western Science} 的選段（Text 1b）中，他闡明了柏拉圖兩個領域的觀念及其與科學的關係。

\noindent 第二篇文本（Text 2）同樣選自 \textit{The Beginnings of Western Science}。它介紹亞里士多德物理學，並以簡明的歷史敘述，說明中世紀科學家後來如何發展運動概念，從而為牛頓物理學的出現奠定基礎。

\noindent Text 3a 選自 \textit{The Birth of a New Physics}。這本書由國際知名的牛頓研究者 I. Bernard Cohen 面向一般讀者撰寫。在這一選段中，Cohen 闡述了牛頓物理學的背景，以及偉大著作 \textit{The Principia: Mathematical Principles of Natural Philosophy}（《自然哲學的數學原理》）背後的主要思想；該書長久以來被視為人類歷史上最重要的著作之一。Text 3b 是這部巨著開首數頁；牛頓在其中建構了一個關於世界的數學體系。讀者將能體會牛頓如何努力為質量、力、運動等抽象概念作出精確定義，也能看到他試圖在數學與哲學之間劃定界線。

